\section{Scoring}

\benchmarkname{} uses deterministic scorers for primary results. The primary
score is answer correctness: did the model give the objectively right answer,
even if it included surrounding prose? The scorer also checks format
correctness and strict correctness as secondary interface-compliance
diagnostics. This keeps the benchmark focused on obvious task failures rather
than making exact response style the headline target.

The paper reports:

\begin{itemize}
  \item answer correctness as the primary ranking metric,
  \item obvious failure rate as one minus answer correctness,
  \item Wilson 95\% confidence intervals for answer correctness,
  \item format accuracy as a secondary compliance diagnostic,
  \item strict accuracy as a secondary exact-interface diagnostic,
  \item provider error count,
  \item estimated cost,
  \item tokens per correct answer, and
  \item cost per correct answer.
\end{itemize}

The current readiness gate requires each scorer used by the paper split to
have at least 20 gold examples. Current scorer-gold coverage and the planned
model panel are listed in the appendix so the main report stays focused on the
benchmark definition and results.

\paragraph{Run protocol.}
Final model sweeps should use a frozen prompt format, frozen model aliases,
recorded run dates, documented generation settings, and a provider-error
accounting policy. After configured retries, final provider errors, refusals,
and timeouts count as incorrect attempts while remaining separately reported.
Smoke runs remain engineering checks and are not ranked against full paper runs.
The current evidence run uses short-answer prompts with a safety-oriented output
cap and records the requested generation settings in the per-model run
manifest. Temperature is fixed where the adapter exposes it so the benchmark is
closer to a deterministic reliability probe than a sampling study.

\paragraph{Thinking controls.}
The report distinguishes configured thinking settings from reported reasoning
tokens. OpenAI reasoning effort, Gemini thinking levels or budgets, Anthropic
thinking budgets, xAI reasoning effort, and OpenRouter-normalized reasoning
settings are request-time controls when explicitly configured. Reported
reasoning tokens are response-time telemetry. They are useful for cost and
efficiency analysis, but they do not by themselves prove which thinking depth
was requested.
